\clearpage
\section{Torque by rotating sphere}
\subsection{Dimensional Analysis}
\(R\) has dimension of length, \(\Omega \) of 1 over time, and \(\mu\) of mass over (length times time). Torque has dimensions of mass times length squared divided by time squared. therefore 
\begin{equation}
    \tau = A\mu\Omega R^3
\end{equation}
\subsection{Symmetry}
The only preferred direction in this problem is generated by the rotation of the sphere. If we move to a rotating frame of reference whose rotation is centered on the center of the sphere, it would appear as if it's in perfect spherical symmetry all around. Therefore, the only direction the velocity of the fluid around the sphere could have is the \(\phi\) direction. But, even then we have symmetry for rotation around \(\phi\), meaning the velocity can't actually depend on \(\phi\). In total
\begin{equation}
    \vb{v} = v(r,\theta)\vu{\phi}
\end{equation}
\subsection{Solution}
Since we're in low reynolds, we'll use the stokes equation 
\begin{equation}
    \grad{p} = \mu\laplacian{\vb{v}}
\end{equation}
rewriting it as 
\begin{equation}
    \laplacian{v_\phi} - \frac{v_\phi}{r^2\sin^2\theta}=0
\end{equation}
now applying Separation of Variables
\begin{equation}
    v_\phi = R(r)\Theta(\theta)
\end{equation}
so 
\begin{align}
    \laplacian{v_\phi} &= \frac{1}{r^2}\del{r}\ins{r^2\del{r}\ins{R\Theta}}+\frac{1}{r^2\sin\theta}\del{\theta}\ins{\sin\theta\del{\theta}\ins{R\Theta}}\\
    &= \frac{1}{r^2}\del{r}\ins{r^2R'\Theta}+\frac{1}{r^2\sin\theta}\del{\theta}\ins{\sin\theta R\Theta'}\\
    &= \frac{\Theta}{r^2}\ins{2rR'+r^2R''}+\frac{R}{r^2\sin\theta}\ins{\cos\theta \Theta'+\sin\theta \Theta''}
\end{align}
so 
\begin{align}
    \frac{R\Theta}{r^2\sin^2\theta}=\frac{\Theta}{r^2}\ins{2rR'+r^2R''}+\frac{R}{r^2\sin\theta}\ins{\cos\theta \Theta'+\sin\theta \Theta''}
\end{align}
or
\begin{align}
    \frac{1}{r^2\sin^2\theta}=\frac{1}{Rr^2}\ins{2rR'+r^2R''}+\frac{1}{\Theta r^2\sin\theta}\ins{\cos\theta \Theta'+\sin\theta \Theta''}
\end{align}
or 
\begin{align}
    \frac{1}{\sin^2\theta}=\frac{1}{R}\ins{2rR'+r^2R''}+\frac{1}{\Theta \sin\theta}\ins{\cos\theta \Theta'+\sin\theta \Theta''}
\end{align}
or
\begin{align}
    \frac{1}{\sin^2\theta}-\frac{\cos\theta \Theta'+\sin\theta \Theta''}{\Theta \sin\theta}=\frac{2rR'+r^2R''}{R}
\end{align}
or 
\begin{align}
    -\frac{\sin\theta\cos\theta \Theta'+\sin^2\theta \Theta''+\Theta}{\Theta \sin^2\theta}=\frac{2rR'+r^2R''}{R} = -\lambda^2
\end{align}
now we can use the separation solution process